\chapter{Introduction}
\label{chap:introduction}

The increasing digitalisation of society is transforming the way information is produced, exchanged, and used \cite{oecd_progress_2026, oecd_scaling_2026}. The growing connectivity between people, devices, and digital services, together with the availability of increasingly large volumes of data and computational resources, has enabled the development of technologies capable of supporting tasks that were previously difficult or impossible to perform \cite{oecd_collective_2024}. Among these technologies, Artificial Intelligence (AI) has emerged as one of the principal drivers of contemporary digital transformation, with applications spanning domains such as healthcare, transportation, finance, manufacturing, and scientific research \cite{hamon_robustness_2020, oecd_progress_2026, annoni_artificial_2018}.

However, the availability of large volumes of data does not automatically translate into reliable information or effective AI systems \cite{peters_dikw_2025}. Data originating from different organisations, applications, devices, and processes may differ in structure, quality, and semantics \cite{oecd_interoperability_2026}. Consequently, the ability to access, integrate, interpret, and reuse data becomes a fundamental prerequisite for exploiting their value \cite{chamola_review_2023}.

This challenge is particularly relevant for AI systems because contemporary Machine Learning (ML) approaches derive computational models from data rather than relying exclusively on explicitly encoded rules \cite{chamola_review_2023}. The resulting systems can achieve remarkable predictive performance, but their behaviour is inherently dependent on the characteristics of the data used throughout their lifecycle \cite{oecd_collective_2024}. Consequently, issues concerning data representation and governance are not merely peripheral concerns but can directly affect the behaviour and reliability of the resulting AI systems \cite{oecd_collective_2024, zha_data-centric_2025, oecd_scaling_2026}.

Moreover, to achieve the effective adoption of AI in society, humans must remain at the centre of this technological evolution. Cutting-edge solutions should support human autonomy and establish a relationship of trust \cite{kaur_trustworthy_2022}. Philosophically, the National Institute of Standards and Technology defines trust as \emph{“the confidence one element has in another, that second element will behave as expected”} \cite{boyens_supply_2015}.

The challenge therefore extends beyond developing increasingly sophisticated AI models \cite{zha_data-centric_2025, kaur_trustworthy_2022}. This thesis investigates the role of data-centric engineering in the development of Trustworthy AI systems. In particular, it explores semantic approaches to transform heterogeneous data into actionable and reliable knowledge, supporting well-informed and proactive decision-making. To demonstrate the contributions of this research, healthcare is adopted as the primary application context. Healthcare represents a particularly relevant setting for this investigation, as it is a fundamental component of civil society in which technological advancements must remain reliable and firmly oriented towards serving individuals and promoting their well-being.

\section{Data-Centricity in the Age of AI}

AI is expected to profoundly reshape the social and economic landscape of contemporary societies, while stimulating innovation and enabling new solutions across a wide range of domains, including medicine, transportation, finance, and machine translation, among others~\cite{annoni_artificial_2018}. AI encompasses a broad family of computational approaches designed to enable machines to perform tasks that typically require cognitive capabilities, such as perception, prediction, reasoning, learning, and decision-making~\cite{hamon_robustness_2020}. In recent years, advances in AI have been driven particularly by machine learning (ML), which enables computational models to infer patterns and relationships from data rather than relying exclusively on manually specified rules~\cite{chamola_review_2023}. The generality of these approaches allows similar underlying principles to be applied to substantially different problems, provided that suitable data are available~\cite{annoni_artificial_2018}.

\textcolor{blue}{TODO: a bridge paragraph that justify the data-centric AI paradigm}

These considerations motivate a shift in attention from the exclusive optimisation of models towards the systematic engineering of the data on which AI systems depend~\cite{zha_data-centric_2025}. This shift has given rise to the concept of \emph{Data-Centric AI}, which places data at the centre of the AI development lifecycle. It is important, however, to distinguish \emph{data-centric} from \emph{data-driven}. The latter primarily refers to the use of data as the basis for developing computational models, whereas a data-centric approach goes further by treating data themselves as an engineering artefact whose quality, structure, representation, documentation, and lifecycle require systematic attention~\cite{zha_data-centric_2025}. This perspective does not imply a replacement of model-centric approaches. Rather, data-centric and model-centric practices are complementary and jointly contribute to the development of reliable and effective AI systems.

%Besides training data, well-designed inference data, the input used in the inference phase of AI systems for making predictions, has facilitated the initial recognition of numerous critical issues in AI and unlocked new model capabilities \cite{zha_data-centric_2025}.

As the volume, variety, and velocity of available data continue to increase, manual approaches alone become increasingly difficult to sustain. Consequently, data-centric AI also requires automated methods to support and streamline the data lifecycle~\cite{zha_data-centric_2025}. This latter is divided in multiple stages, including \emph{Training Data Development} and \emph{Data Maintenance}~\cite{zha_data-centric_2025}.

\paragraph{Training Data Development} concerns the collection, integration, and preparation of data for AI model development~\cite{zha_data-centric_2025}. This process can require substantial engineering effort; a frequently cited estimate suggests that it accounts for approximately 80\% of data scientists' work~\cite{press_cleaning_2020}. It comprises (a) \emph{Data collection}, which gathers data from relevant sources and requires domain knowledge to determine what should be collected and how it should be represented; (b) \emph{Data integration}, which combines data from multiple sources into a unified representation by aligning attributes and transforming source-specific values; and (c) \emph{Data preparation}, which cleans and transforms the resulting data into a form suitable for model development.

\paragraph{Data Maintenance} recognises that data are continuously generated, updated, and reused in operational environments~\cite{zha_data-centric_2025}. It therefore comprises (a) \emph{Data understanding}, which supports the interpretation and assessment of complex datasets through activities such as summarisation and visualisation. It enables stakeholders to evaluate data characteristics, provenance, limitations, and suitability for their intended use; (b) \emph{Data storage and retrieval}, which provide the scalable infrastructure required to make data available to AI systems as data volumes and generation rates increase; and (c) \emph{Data quality assurance}, which involves the continuous assessment and improvement of data quality. It may consider objective dimensions, such as accuracy, timeliness, consistency, and completeness, as well as human-centred dimensions, including trustworthiness, understandability, and accessibility.


%\paragraph{Data understanding} is equally important because real-world datasets are often large, heterogeneous, and difficult to interpret. A comprehensive understanding of the characteristics, provenance, limitations, and intended use of data is necessary to assess their suitability for a particular task and to identify potential sources of error or bias~\cite{zha_data-centric_2025}. Such understanding should extend beyond the data themselves to the contexts in which they were collected, processed, and intended to be used. This is particularly important when data are subsequently employed in environments that differ from those in which they were originally collected.

%\paragraph{Data quality assurance} concerns the continuous assessment and maintenance of data quality throughout the lifecycle of an AI system. In real-world settings, both data sources and the infrastructures used to process them are subject to frequent and continuous change. Consequently, it is not sufficient to produce high-quality training or inference data at a single point in time; their quality must be monitored and maintained as the underlying data and operational environment evolve~\cite{zha_data-centric_2025}. This requires appropriate metrics and evaluation procedures for characterising different dimensions of data quality. Importantly, data quality cannot be understood solely through technical measures. It also encompasses human-centred dimensions, including trustworthiness, understandability, and accessibility~\cite{zha_data-centric_2025}. Trustworthiness concerns the reliability and accuracy of information provided by a data source; understandability concerns the extent to which users can comprehend the collected data; and accessibility concerns the ability of relevant users to access and effectively use the data.

Taken together, these activities illustrate that data in AI systems are not merely passive inputs to computational models. They are artefacts that must be acquired, transformed, documented, validated, monitored, maintained, and interpreted throughout their lifecycle. Data management is therefore a fundamental component of trustworthy AI, motivating principles and practices for the appropriate storage, management, and use of data~\cite{hamon_robustness_2020}. Building on these principles, data and AI governance are gaining increasing prominence as mechanisms for ensuring that the design, development, deployment, and ongoing management of AI systems remain aligned with the values, responsibilities, and expectations of organisations and society~\cite{hamon_robustness_2020}.

\section{From Data to AI Governance}

The central role of data in contemporary AI systems makes their management an essential component of responsible AI development. This has motivated the design of principles and practices for \emph{Data Governance}, aimed at ensuring that data are appropriately managed throughout their lifecycle according to requirements concerning quality, accountability, security, privacy, and lawful use \cite{hamon_robustness_2020}. Data governance, therefore, constitutes an important foundation upon which broader AI governance and trustworthiness can be established \cite{hamon_robustness_2020}.

Data governance is particularly relevant when AI systems process personal or otherwise sensitive information. In the European context, the General Data Protection Regulation (GDPR) established a comprehensive framework for the processing of personal data and the protection of the rights of data subjects \cite{data_protection_working_party_article29_nodate}. Its principles are particularly relevant to AI systems that support or automate decisions concerning individuals, where questions of transparency, accountability, and human involvement become increasingly important \cite{panigutti_role_2023}. Specifically, when an automated decision is taken, and its outcomes may have a negative consequences for the data subject, the model behind it must be transparent and the decision interpretable\cite{panigutti_role_2023}. Moreover, freedoms of data subjects must be safeguarded by granting their right to obtain the explanation\cite{kaminski_right_2019} and contest the decision taken\cite{data_protection_working_party_article29_nodate}. 

The development of AI systems introduces an additional layer of governance. While Data Governance focuses primarily on how data are collected, managed, protected, and used, \emph{AI Governance} extends these concerns to the systems that process the data and to the decisions and actions resulting from their operation\cite{hamon_robustness_2020}.  This transition is reflected in the European Union's development of  AI Act that prescribes high-risk AI systems to be developed in such a way that they enable users to interpret their output correctly and use them appropriately\cite{corbucci_semantic_2023}.  

Through the EU AI Act, the European Union aims to ensure that fundamental rights, democracy, the rule of law and environmental sustainability are protected from high-risk AI\cite{oecd_collective_2024}. The Act categorises AI systems based on risk levels, imposing bans for systems that present unacceptable risks\cite{oecd_collective_2024}. Moreover, the AI Act adopts concepts commonly used in the eXplainable AI (XAI) literature, such as transparency, opacity, and comprehensibility \cite{panigutti_role_2023}. However, this overlap is primarily terminological: the AI Act does not prohibit the use of black-box AI systems, nor does it prescribe the use of specific XAI tools or techniques \cite{panigutti_role_2023}. The AI act frames transparency primarily as an obligation to provide users with appropriate information and documentation\cite{panigutti_co-design_2023}. These requirements demonstrate that trustworthiness cannot be achieved through model performance alone, but requires coordinated attention to the data, the technical system, its documentation, and the humans involved in its operation\cite{panigutti_role_2023,oecd_collective_2024}. 

\section{Trustworthy AI}\label{sec:trustworthy-ai}

As anticipated in the previous section, the objective is not only to ensure that the underlying data are appropriately managed, but also that AI systems operate in accordance with legal requirements, ethical principles, technical constraints, and societal expectations \cite{hamon_robustness_2020}. This shift from Data Governance to AI Governance is reflected in the increasing international and European attention devoted to \emph{Trustworthy AI}. The need for trustworthy systems arises from the observation that technical performance alone is insufficient to determine whether an AI system is appropriate for deployment\cite{kaur_trustworthy_2022}. An AI system may achieve high predictive accuracy yet remain unsuitable for a specific context, such as healthcare, due to biased data, inadequate privacy protection, limited robustness, lack of transparency, or insufficient mechanisms for human intervention
\cite{kaur_trustworthy_2022, oecd_progress_2026, li_trustworthy_2023}.

Trustworthy AI therefore extends the assessment of AI systems beyond predictive performance. Several organizations and research communities have proposed principles and frameworks for trustworthy AI\cite{european_commission_ethics_nodate,lekadir_future-ai_2025, oecd_oecd_2019,unesco_recommendation_2022}, with an emerging convergence around dimensions including transparency and explainability \cite{kaur_trustworthy_2022}. From this perspective, trustworthiness is not a single technical property that can be added to an AI system after its development. Rather, it results from the interaction of multiple properties that must be considered throughout the AI lifecycle.

\paragraph{The EU framework}
The European Commission's Ethics guidelines for Trustworthy AI is a comprehensive framework for the development, deployment, and use of trustworthy AI systems \cite{european_commission_ethics_nodate}. According to these guidelines, AI systems should be (a) \emph{lawful}, by respecting applicable laws and regulations; (b) \emph{ethical}, by adhering to ethical principles and values; and (c) \emph{robust}, both from a technical perspective and with regard to their social environment.

To operationalize these principles, the guidelines identify seven key requirements for Trustworthy AI \cite{european_commission_ethics_nodate,kaur_trustworthy_2022}: (1) \emph{human agency and oversight}, ensuring that AI systems support human autonomy and allow for appropriate human intervention; (2) \emph{technical robustness and safety}, ensuring reliable and safe operation, including under unexpected conditions or external disturbances; (3) \emph{privacy and data governance}, ensuring the protection of personal data and compliance with applicable data protection requirements; (4) \emph{transparency}, providing relevant information about the system's capabilities, limitations, and operation to enable stakeholders to appropriately understand and use its outputs; (5) \emph{diversity, non-discrimination, and fairness}, preventing unjustified discrimination and accounting for the diversity of individuals and groups affected by AI systems; (6) \emph{societal and environmental well-being}, ensuring that AI systems contribute positively to society while avoiding disproportionate societal or environmental harms; and (7) \emph{accountability}, establishing appropriate mechanisms for assigning responsibility for the development, deployment, operation, and outcomes of AI systems.

\paragraph{The role of data on trust}
The seven requirements of the EU framework highlight that trustworthiness is a system-level property. At the same time, several of these requirements are directly or indirectly shaped by how data are collected, represented, transformed, documented, and maintained \cite{kaur_trustworthy_2022}. Privacy and data governance explicitly concern the management of data, but their influence extends beyond this requirement. Technical robustness depends, among other factors, on the quality and suitability of the data used for training and operation. Fairness can be compromised by biases and under-representation in datasets, while transparency and accountability can be strengthened when the provenance, meaning, transformations, and intended uses of data can be established and communicated. Similarly, effective human oversight relies on the availability of sufficient contextual information to enable users to appropriately assess and interpret AI outputs \cite{guidotti_survey_2019}.

Building on these observations, this thesis positions data-centricity as a key enabler of Trustworthy AI; rather, the central argument is that trustworthiness cannot be reliably built on data that are poorly understood or inaccessible. This perspective calls for a wider interpretation of the data-centric paradigm. Data engineering should not be limited to improving the statistical properties of datasets used for model training. Still, it should also address how data are represented and transformed into knowledge that supports human and machine reasoning.

\paragraph{Explainable AI}
Among the requirements of Trustworthy AI, transparency is particularly relevant to the relationship between AI systems and their human users. The question of how computational systems can make their reasoning understandable to humans is not new \cite{hamon_robustness_2020}. Recently, the field of XAI has expanded to encompass technical approaches for designing interpretable models and generating explanations \cite{guidotti_survey_2019}, methodologies for evaluating explanations \cite{lekadir_future-ai_2025}, and contributions from social sciences concerning the nature and role of explanations \cite{miller_explanation_2019}. XAI is therefore inherently interdisciplinary, involving computer science, human-computer interaction, social sciences, law, and philosophy \cite{panigutti_role_2023}.

A central difficulty in XAI is that there is no universally accepted formal definition of an \emph{explanation}. Although explanation has been studied extensively in disciplines such as cognitive science and social psychology \cite{halpern_causes_2005}, many technical approaches developed within the AI community do not explicitly rely on a formal theory of explanation \cite{miller_explanation_2019}. Instead, they generally aim to provide information about the behaviour of an AI system in a form that can be understood by humans\cite{chamola_review_2023, corbucci_semantic_2023, johannssen_explainable_2025}. Consequently, what constitutes an adequate explanation depends on the application, the intended user, and the purpose for which the explanation is required \cite{murdoch_interpretable_2019}. 

\section{AI in Healthcare}
\label{ai_in_healthcare}

The healthcare sector represents a particularly relevant context for investigating the opportunities and challenges associated with AI. Healthcare is not only a fundamental component of modern societies but also a complex socio-technical ecosystem involving diverse professions, processes, and practices that rely on the production and exchange of large volumes of information. Consequently, technological progress in healthcare cannot be evaluated solely in terms of computational performance: the quality of care is closely connected to the quality, availability, and interoperability of the information on which clinical and organizational processes depend. AI systems must operate within this complex environment, where their outputs can have significant consequences for human health and well-being. These characteristics make healthcare a particularly demanding and consequential setting for investigating the data-centric foundations of Trustworthy AI.

\paragraph{Some statistics} 
Healthcare is undergoing a profound digital transformation driven by the increasing availability of health data and the adoption of technologies such as AI, the Internet of Things (IoT), wearable devices\cite{oecd_progress_2026}. Clinical information already accounts for approximately 30\% of global data, and its volume continues to grow rapidly \cite{thomason_data_2024}. The healthcare sector itself represents a substantial economic domain, with global expenditure estimated at approximately USD 8.3 trillion\cite{thomason_data_2024}. At the same time, IoT technologies are expected to generate up to USD 1.8 trillion in economic value by 2030\cite{seixas-lopes_musculoskeletal_2024}, while the wearable medical device market, valued at USD 14.4 billion in 2020, has been projected to grow at an annual rate of 18.9\% \cite{thomason_data_2024}. Finally, it is projected that the economic value of the AI market will grow USD 183 billion in 2023 to USD 995 billion in 2028, which represents more than 450\% growth\cite{oecd_collective_2024}.

The relevance of this transformation is further reinforced by the structural pressures faced by healthcare systems. Across the European Union, healthcare expenditure represented approximately 10.4\% of Gross Domestic Product (GDP) in 2022, while healthcare spending increased by approximately 39.9\% between 2014 and 2022 \cite{oecd_progress_2026}. At the same time, demographic ageing is increasing the demand for healthcare services: people aged 65 years and over represented approximately 21\% of the European population in 2023 and are projected to account for approximately 29\% by 2050 \cite{oecd_progress_2026}. These demographic and economic trends are accompanied by workforce pressures, with an OECD projection estimating a potential shortage of approximately 3.5 million health professionals by 2030 \cite{oecd_progress_2026}. In this context, AI is increasingly considered not merely as a technological innovation, but as an enabler for new models of care.

\paragraph{AI opportunities}

The potential impact of AI in healthcare extends across the entire care pathway \cite{oecd_progress_2026}. Estimates suggest that AI could automate or substantially support activities corresponding to up to 36\% of tasks performed in health and social care \cite{kumar_chebrolu_smart_2020}. Such applications could reduce the burden of repetitive administrative and analytical activities, allowing healthcare professionals to devote more attention to patient-centred tasks \cite{oecd_progress_2026, hamon_robustness_2020}. Beyond task automation, the integration of predictive models with technologies such as wearable devices, remote monitoring, and Digital Twins may support more personalized and preventive approaches to care \cite{oecd_progress_2026}.

AI may also contribute directly to clinical decision support. Recent assessments suggest that AI could address approximately 20\% of clinical needs in the coming years\cite{gonzalez-gonzalo_trustworthy_2022}. Further, AI can reduce medical errors due to miscommunication; estimated at 30\% of all errors \cite{european_alliance_for_access_to_safe_medicines_medication_2022}. Moreover, AI can unlock value from the 97\% of the health data assets that are unused for decision making \cite{oecd_collective_2024}. This gap highlights a central paradox of digital healthcare: \emph{the availability of data is increasing rapidly, while the capacity to transform those data into usable information and actionable knowledge remains comparatively limited}.

Examples of potential applications include (1) \emph{empowering individuals and personalising care} \cite{oecd_collective_2024}: AI can provide insights into personal health data, facilitate understanding of medical records and healthcare systems, and support informed care decisions \cite{marfoglia_towards_2025,krumholz_big_2014,topol_high-performance_2019}, potentially improving health outcomes, reducing costs, and increasing patient satisfaction \cite{hibbard_what_2013}; (2) \emph{streamlining processes and enhancing decision-making}: AI can reduce clinicians' administrative burden, potentially freeing 10--30\% of their time for patient care, and support evidence-based practice by identifying relevant scientific literature \cite{shanafelt_potential_2016,oecd_collective_2024}; and (3) \emph{advancing drug discovery}: AI can accelerate the identification of promising drug candidates while reducing the resources required for early-stage discovery, as demonstrated by the identification of 240 potential antibiotics from over 7,000 candidates in approximately two hours \cite{liu_deep_2023}.